\documentclass[11pt,a4paper]{article}

\usepackage[utf8]{inputenc}
\usepackage[T1]{fontenc}
\usepackage[french]{babel}
\usepackage{lmodern}
\usepackage[margin=2cm]{geometry}
\usepackage{enumitem}
\usepackage{xcolor}
\definecolor{lienbleu}{RGB}{0,70,140}
\usepackage{xurl}
\usepackage[colorlinks=true, urlcolor=lienbleu, linkcolor=lienbleu, citecolor=lienbleu]{hyperref}

% Publications : DOI cliquable et nom de l'auteur en gras
\newcommand{\doi}[1]{\href{https://doi.org/#1}{doi:\nolinkurl{#1}}}
\newcommand{\me}{\textbf{V.~De~Bono}}
\newenvironment{publis}[1]%
  {\begin{itemize}[label={}, leftmargin=1.2em, labelwidth=0pt, labelsep=0pt,
     itemindent=-1.2em, itemsep=0.45em, topsep=0.3em]}%
  {\end{itemize}}

% Mise en page 
\setlist[description]{leftmargin=2.8cm, labelwidth=2.5cm, labelsep=0.3cm,
  align=left, font=\normalfont\bfseries, itemsep=0.5em}
\setlist[itemize]{topsep=0.2em, itemsep=0em}
\setlength{\parindent}{0pt}
\pagestyle{empty}

\title{CV}
\author{Victor De Bono}
\date{2026}

\begin{document}

\begin{center}
  {\LARGE\textbf{Victor De Bono}}\\[0.6em]
\end{center}

% ---------------------------------
\textbf{Maître de conférences en sciences et techniques pour l'architecture (STA)}, ENSAP Lille\\
\textbf{Chercheur}, LACTH : Laboratoire d'architecture, conception, territoire, histoire, matérialité\\
\textbf{Chercheur associé}, UMR MAP-ARIA : Modèles et simulations pour l'architecture et le patrimoine\\
\textbf{ORCID} \href{https://orcid.org/0009-0006-7596-9395}{0009-0006-7596-9395}


% ---------------------------------
\subsection*{Articles de revues scientifiques}

\begin{publis}{A}
  \item \me, N. Ducoulombier, R. Mesnil, J.-F. Caron.
  A roadmap for the structural analysis and justification of thin anisotropic 3D printed reinforced concrete shells.
  Preprint, \emph{SSRN}, 5204311.
  \href{https://papers.ssrn.com/sol3/papers.cfm?abstract_id=5204311}{ssrn.com/abstract=5204311}

  \item S. Phetsinorath, J. Archez, N. Ducoulombier, \me, L. Demont, \emph{et al.} (2026).
  A 2K strategy for 3D printing of metakaolin-based geopolymers.
  \emph{Materials and Structures}, 59(5), 272.
  \href{https://scholar.google.com/scholar?q=A+2K+strategy+for+3D+printing+of+metakaolin-based+geopolymers}{Google Scholar}

  \item L. Caneda-Martínez, M. Hassan, L. Demont, E. Keita, P. Belin, \me, R. Mesnil, J.-F. Caron, B. González-Fonteboa, N. Roussel (2026).
  Fast penetration testing of printable concretes with a portable device: robustness and calibration.
  \emph{Cement and Concrete Research}, 202, 108141.
  \doi{10.1016/j.cemconres.2026.108141}

  \item \me, R. Mesnil, N. Ducoulombier, J.-F. Caron (2025).
  Design, fabrication and testing of a 3D printed concrete shell with continuous fiber-reinforcement.
  Preprint, \emph{SSRN}, 6213618.
  \href{https://papers.ssrn.com/sol3/papers.cfm?abstract_id=6213618}{ssrn.com/abstract=6213618}

  \item R. Wolfs, J. Versteege, M. Santhanam, S. Bhattacherjee, F. Bos, S. Muthukrishnan, C. Menna, O. Ozturk, N. Özyurt, J. Roupec, C. Richter, J. Jungwirth, L. de Miranda, R. Ammann, J.-F. Caron, \me, \emph{et al.} (2025).
  Mechanical properties of 3D printed concrete: a RILEM TC 304-ADC interlaboratory study - flexural and tensile strength.
  \emph{Materials and Structures}, 58(5), 182.
  \doi{10.1617/s11527-025-02687-w}

  \item V. Mechtcherine, S. Muthukrishnan, A. Robens-Radermacher, R. Wolfs, J. Versteege, C. Menna, O. Ozturk, N. Özyurt, J. Roupec, C. Richter, J. Jungwirth, L. de Miranda, R. Ammann, J.-F. Caron, \me, \emph{et al.} (2025).
  Mechanical properties of 3D printed concrete: a RILEM TC 304-ADC interlaboratory study - compressive strength and modulus of elasticity.
  \emph{Materials and Structures}, 58(5), 181.
  \doi{10.1617/s11527-025-02688-9}

  \item A. Perrot, Y. Jacquet, J.-F. Caron, R. Mesnil, N. Ducoulombier, \me, J. Sanjayan, S. Ramakrishnan, H. Kloft, J. Gosslar, S. Muthukrishnan, V. Mechtcherine, T. Wangler, J. L. Provis, K. Dörfler, E. Krakovská, N. Roussel, E. Keita (2024).
  Snapshot on 3D printing with alternative binders and materials: earth, geopolymers, gypsum and low carbon concrete.
  \emph{Cement and Concrete Research}, 185, 107651.
  \doi{10.1016/j.cemconres.2024.107651}

  \item \me, N. Ducoulombier, R. Mesnil, J.-F. Caron (2024).
  Methodology for formulating low-carbon printable mortar through particles packing optimization.
  \emph{Cement and Concrete Research}, 176, 107403.
  \doi{10.1016/j.cemconres.2023.107403}
\end{publis}

% ---------------------------------
\subsection*{Chapitres d'ouvrage}

\begin{publis}{O}
  \item N. Ducoulombier, \me, F. Z. Kachkouch, Y. Jacquet, A. Perrot (2025).
  From laboratory to practice: characterizing fresh and cured printed materials.
  In A. Perrot, Y. Jacquet (eds.), \emph{3D Concrete Printing: State of the Art and Applications}, chap.~3, pp.~67-101. ISTE/Wiley.
  \doi{10.1002/9781394352081.ch3}

\end{publis}

% ---------------------------------
\subsection*{Conférences}

\begin{publis}{C}
  \item A. Sepidbar, M. Pierre, \me, N. Ducoulombier, R. Mesnil, \emph{et al.} (2026).
  An in-line method for estimation of early-age stiffness of 3D concrete printing.
  \emph{RILEM International Conference on Concrete and Digital Fabrication}, pp.~46-53.
  \href{https://link.springer.com/chapter/10.1007/978-3-032-36738-9_6}{DC2026 Springer}

  \item \me, R. Mesnil, J.-F. Caron, N. Ducoulombier (2026).
  Conception, fabrication et essais mécaniques sur une coque mince imprimée 3D en béton fibré.
  \emph{SCAN'26}, ENSA Montpellier. \emph{SHS Web of Conferences}, 241, 02007.
  \href{https://www.shs-conferences.org/articles/shsconf/abs/2026/15/shsconf_scan26_02007/shsconf_scan26_02007.html}{shs-conferences.org}

  \item J.-F. Caron, N. Ducoulombier, L. Demont, \me, R. Mesnil (2026).
  Using 3D printing for a new cementitious composite and structural design.
  \emph{Proceedings of the 12th International Conference on FRP Composites in Civil Engineering (CICE)}, vol.~777, pp.~1615-1623. Springer.
  \doi{10.1007/978-3-032-09399-8_152}

  \item \me, N. Ducoulombier, L. Demont, R. Mesnil, J.-F. Caron (2025).
  Coques cimentaires imprimées en 3D et renforcées par fibres longues.
  \emph{Comptes rendus des Journées nationales sur les composites (JNC 2025)}, Gif-sur-Yvette.
  \href{https://jnc2025.sciencesconf.org/602015/document}{jnc2025.sciencesconf.org/602015}

  \item \me, N. Ducoulombier, S. Loulha, R. Mesnil, J.-F. Caron (2025).
  Toward 3D printable low-carbon mortar: method and application.
  \emph{Proceedings of the RILEM Spring Convention and Conference 2024}, Milan. RILEM Bookseries, 56, pp.~221-230. Springer.
  \doi{10.1007/978-3-031-70281-5_25}

  \item L. Demont, \me, N. Ducoulombier, R. Mesnil, J.-F. Caron (2024).
  Anisotropic concrete: a composite look on the engineering and architecture of concrete construction.
  \emph{Built Environment Additive Manufacturing (BE-AM) Symposium 2024}, N. Gaudillière-Jami, A. Wolf (eds.), pp.~94-103.

  \item \me, L. Demont, N. Ducoulombier, R. Mesnil, J.-F. Caron (2024).
  Large-scale 3D printing of a shell structure reinforced by the flow-based pultrusion process.
  \emph{Supplementary Proceedings of the 4th RILEM International Conference on Concrete and Digital Fabrication (DC 2024)}, Munich.
  \doi{10.24355/dbbs.084-202408201037-0}

  \item \me, N. Ducoulombier, R. Mesnil, J.-F. Caron (2024).
  Tailored low-carbon footprint cementitious material for 3D printing: strategies for rheology adjustment and mechanical performance estimation.
  \emph{Supplementary Proceedings of the 4th RILEM International Conference on Concrete and Digital Fabrication (DC 2024)}, Munich.
  \doi{10.24355/dbbs.084-202408130656-0}

  \item L. Caneda-Martínez, M. Hassan, L. Demont, E. Keita, P. Belin, \me, \emph{et al.} (2024).
  Optimization of fast penetration tests for rheology control in 3D printing concrete.
  \emph{Supplementary Proceedings of the 4th RILEM International Conference on Concrete and Digital Fabrication (DC 2024)}, Munich.
  \href{https://leopard.tu-braunschweig.de/receive/dbbs_mods_00077815}{leopard.tu-braunschweig.de}

  \item \me, R. Mesnil, N. Ducoulombier, J.-F. Caron (2024).
  Design, structural optimization and fabrication of concrete shells through fibre-reinforced 3D printing.
  \emph{Proceedings of the IASS Annual Symposium 2024: Redefining the Art of Structural Design}, ETH Zurich, n°~4, pp.~1-10.
  \href{https://www.ingentaconnect.com/contentone/iass/piass/2024/00002024/00000004/art00009}{ingentaconnect.com}

  \item J.-F. Caron, N. Ducoulombier, L. Demont, \me, R. Mesnil (2024).
  Reinvent reinforced concrete with robotics and 3D printing.
  \emph{Proceedings of the 41st International Symposium on Automation and Robotics in Construction (ISARC 2024)}, Lille.
  \doi{10.22260/ISARC2024/0021}

  \item J.-F. Caron, N. Ducoulombier, L. Demont, \me, R. Mesnil (2023).
  3D printing of continuous-fibers cementitious composites: anisotropic 3D mortar.
  \emph{Visions and Strategies for Reinforcing Additively Manufactured Constructions}, \emph{Open Conference Proceedings}, 3, pp.~65-74.
  \doi{10.52825/ocp.v3i.193}

  \item \me, N. Ducoulombier, R. Mesnil, J.-F. Caron (2023).
  Formulation de béton pour l'impression 3D.
  \emph{AFJC}, Paris. \emph{Academic Journal of Civil Engineering}, 41(1), pp.~447-456.
  \doi{10.26168/ajce.41.1.4}
\end{publis}

% ---------------------------------
\subsection*{Thèse de doctorat}

\begin{publis}{T}
  \item \me{} (2025).
  Étude des procédés d'impression 3D de bétons fibrés pour la construction. Formulation, renforcement et application.
  Thèse de doctorat, École nationale des ponts et chaussées, laboratoire Navier. Dir. J.-F. Caron, co-encadrement N. Ducoulombier et R. Mesnil.
  \href{https://theses.hal.science/tel-05329680/}{theses.hal.science/tel-05329680}
\end{publis}


\section*{Enseignement}

\begin{description}
  \item[2026 - auj.] \textbf{Maître de conférences titulaire STA-CIMA}, ENSAP Lille
  \begin{itemize}
    \item TD Franchissement, licence 1 - 36 h TD
    \item CM et TD Thermique du bâtiment, licence 2 - 22 h CM, 24 h TD
    \item CM et TD Système constructif, licence 2 - 22 h CM, 24h TD
    \item TD Architecture bioclimatique, licence 3 - 12 h TD
    \item TD Enveloppes, licence 3 - 24 h TD
    \item CM Analyse de cycle de vie, master 1 - 22 h CM
    \item Exploratoire Fabrication additive, master 1 - 48 h TD
  \end{itemize}

  \item[2025 - 2026] \textbf{Maître de conférences associé STA}, ENSA Bretagne
  \begin{itemize}
    \item TD Cultures matérielles et constructives, licence 1 - 24 h TD
    \item TD Modélisation BIM, Revit, licences 2 et 3 - 96 h TD
    \item CM Thermique du bâtiment, licence 2 - 24 h TD
    \item Atelier Architecte-ingénieur, master 1 - 96 h TD
    \item TD Modélisation BIM et paramétrique, Revit/Dynamo, master 1 - 48 h TD
    \item Approfondissement \emph{Hybridations : artefact, calcul, structure}, master 1 - 12 h TD
  \end{itemize}

  \item[2019 - 2025] \textbf{Enseignant vacataire}, ENSA Lyon
  \begin{itemize}
    \item Intensif numérique, modélisation paramétrique (Rhinoceros, Grasshopper), licences 2 et 3.
    \item CM Continuum numérique et cartographie du numérique, licence 2 et master FPC.
  \end{itemize}

  \item[2023 - 2024] \textbf{Enseignant vacataire}, ENSA Lyon
  \begin{itemize}
    \item Accompagnement de projet Archifolies, master 1.
  \end{itemize}

  \item[2022 - 2025] \textbf{Enseignant vacataire}, École nationale des ponts et chaussées
  \begin{itemize}
    \item TP Connaissance des matériaux, 1\textsuperscript{re}, et  2\textsuperscript{e} années.
    \item Workshop impression 3D béton et argile pour la construction, mastère spécialisé \emph{Design by Data}.
  \end{itemize}

  \item[2018 - 2021] \textbf{Enseignant vacataire}, ENTPE
  \begin{itemize}
    \item Outils de modélisation BIM, 3\textsuperscript{e} année.
  \end{itemize}

  \item[2022 - 2024] \textbf{Intervenant}, ÉTS Montréal et Faculté de l'aménagement, Université de Montréal
  \begin{itemize}
    \item CM Fabrication additive et construction : historique, technologies et applications, 2\textsuperscript{e} année.
  \end{itemize}

  \item[2021 - 2022] \textbf{Workshop}, École Camondo
  \begin{itemize}
    \item Outils de conception paramétrique et impression 3D pour l'architecture intérieure, 1\textsuperscript{re} à 3\textsuperscript{e} années.
  \end{itemize}

  \item[2018 - 2020] \textbf{Enseignant vacataire}, Institut d'urbanisme de Lyon, Université Lyon 2
  \begin{itemize}
    \item Outils de représentation pour l'urbanisme (suite Adobe), licence 2.
  \end{itemize}
\end{description}


\section*{Expérience professionnelle}

\begin{description}
  \item[2019 - 2025] \textbf{Chef de projet}, XtreeE
  \begin{itemize}
    \item Gestion de projets d'impression 3D grande échelle pour la construction, de la conception à la livraison.
    \item Conception et modélisation paramétrique appliquées aux éléments architecturaux et structurels, phases APS, APD.
    \item Réalisation de carnets de détails, systèmes d'accroche et d'assemblage, phases PRO et EXE.
    \item Gestion de sessions d'impression 3D grande échelle multimatériaux (béton, plâtre, terre, géopolymère).
    \item Formation à la fabrication robotique.
  \end{itemize}

  \item[2016 - 2019] \textbf{BIM management}, freelance
  \begin{itemize}
    \item Accompagnement et implémentation de processus BIM : chartes BIM pour la MOE, cahiers des charges BIM pour la MOA, suivi des livrables.
    \item Vérification de maquettes numériques IFC (lycées, logements, bureaux), toutes phases de l'APS au DOE, marchés publics et privés.
    \item Production de maquettes interactives VR.
  \end{itemize}

  \item[2015 - 2016] \textbf{BIM manager}, Astus Construction
  \begin{itemize}
    \item Rédaction de chartes BIM pour la MOE et de cahiers des charges BIM pour la MOA.
  \end{itemize}
\end{description}

% ---------------------------------
% ---------------------------------
\section*{Diplômes}

\begin{description}
  \item[2022 - 2025] \textbf{École nationale des ponts et chaussées}, diplôme national de doctorat
  \begin{itemize}
    \item Thèse : \emph{Étude des procédés d'impression 3D de bétons fibrés pour la construction. Formulation, renforcement et application}. Préparée au laboratoire Navier, soutenue le 12 septembre 2025.
  \end{itemize}

  \item[2013 - 2018] \textbf{INSA Lyon}, diplôme d'ingénieur
  \begin{itemize}
    \item Génie civil et urbanisme, spécialité bâtiment, structure et énergie.
  \end{itemize}

  \item[2011 - 2016] \textbf{ENSA Lyon}, diplôme d'État d'architecte
  \begin{itemize}
    \item Spécialité architecture numérique et formes complexes.
    \item Sélectionné au prix de la jeune architecture.
  \end{itemize}
\end{description}

% ---------------------------------


\end{document}